
\documentclass{article}
\usepackage{colortbl}
\usepackage{makecell}
\usepackage{multirow}
\usepackage{supertabular}

\begin{document}

\newcounter{utterance}

\centering \large Interaction Transcript for game `codenames', experiment `opponent\_easy', episode 4 with claude{-}opus{-}4{-}8{-}azure.
\vspace{24pt}

{ \footnotesize  \setcounter{utterance}{1}
\setlength{\tabcolsep}{0pt}
\begin{supertabular}{c@{$\;$}|p{.15\linewidth}@{}p{.15\linewidth}p{.15\linewidth}p{.15\linewidth}p{.15\linewidth}p{.15\linewidth}}
    \# & $\;$A & \multicolumn{4}{c}{Game Master} & $\;\:$B\\
    \hline

    \theutterance \stepcounter{utterance}  
    & & & \multicolumn{2}{p{0.3\linewidth}}{
        \cellcolor[rgb]{0.9,0.9,0.9}{
            \makecell[{{p{\linewidth}}}]{
                \texttt{\tiny{[GM$|$GM]}}
                {'hidden': {'team': ['kulüp', 'devrim', 'berlin', 'antarktika', 'stadyum', 'öğrenci', 'mikroskop', 'savaş', 'bal'], 'innocent': ['tokyo', 'posta', 'taç', 'süper kahraman', 'ardıç kuşu', 'nal', 'çapraz'], 'opponent': ['sıra', 'ayı', 'tahvil', 'bilim adamı', 'bıçak', 'tarih', 'kaçakçı', 'olympos'], 'assassin': ['reçel']}, 'revealed': {'team': {'team': [], 'innocent': [], 'opponent': [], 'assassin': []}, 'opponent': {'team': [], 'innocent': [], 'opponent': [], 'assassin': []}}}
            }
        }
    }
    & & \\ \\

    \theutterance \stepcounter{utterance}  
    & & \multicolumn{4}{p{0.6\linewidth}}{
        \cellcolor[rgb]{0.9,0.9,0.9}{
            \makecell[{{p{\linewidth}}}]{
                \texttt{\tiny{[P1$\langle$GM]}}
                \texttt{Aşağıdaki listeden takımının kelimelerinden bazılarıyla ilişkili tek bir tek kelimelik ipucu ver. İpucunun kaç kelimeyle ve hangi kelimelerle ilişkili olacağını sen seçebilirsin.} \\
\texttt{Amaç, diğer takım kendi kelimelerinin hepsini bulmadan önce, verdiğin ipuçlarına dayanarak takım arkadaşının senin takımına ait tüm kelimeleri tahmin etmesini sağlamak. Takım arkadaşın bir "rakip" ya da bir "dikkat dağıtıcı" kelime tahmin ederse, tahmin sırası biter ve sıra rakip takıma geçer.  Eğer partnerin bir "öldüren" kelime tahmin ederse, oyun biter ve takımın anında kaybeder.} \\
\texttt{İpucun ne kadar çok kelimeyle ilişkili olursa oyunu o kadar hızlı kazanabilirsin, ancak çağrışımların takım arkadaşın için tahmin etmesi o kadar da zor olabilir; bu yüzden ipucu kelimeni akıllıca seç. İpucu kelimesi hedef kelimelerle anlamsal olarak ilişkili olmalıdır; listelerdeki kelimelerden biri olamaz ve onların parçalarını içeremez.} \\
\\ 
\texttt{Tek kelimelik ipucunu ve ilgili hedef kelimelerin virgülle ayrılmış listesini her zaman aşağıdaki biçimde ver ve yanıtlarını olabildiğince kısa tut; bu formda gereken dışında asla başka metin ekleme:} \\
\\ 
\texttt{İPUCU: <KELİME>} \\
\texttt{HEDEFLER: <HEDEFLER>} \\
\\ 
\texttt{Takımının kelimeleri: kulüp, devrim, berlin, antarktika, stadyum, öğrenci, mikroskop, savaş, bal.} \\
\texttt{Rakibinin takım kelimeleri: sıra, ayı, tahvil, bilim adamı, bıçak, tarih, kaçakçı, olympos.} \\
\texttt{Dikkat dağıtıcı kelimeler: tokyo, posta, taç, süper kahraman, ardıç kuşu, nal, çapraz.} \\
\texttt{Öldüren kelimeler: reçel.} \\
            }
        }
    }
    & \\ \\

    \theutterance \stepcounter{utterance}  
    & \multicolumn{4}{p{0.6\linewidth}}{
        \cellcolor[rgb]{0.9,0.9,0.9}{
            \makecell[{{p{\linewidth}}}]{
                \texttt{\tiny{[P1$\rangle$GM]}}
                \texttt{İPUCU: mücadele} \\
\texttt{HEDEFLER: devrim, savaş} \\
            }
        }
    }
    & & \\ \\

    \theutterance \stepcounter{utterance}  
    & & & \multicolumn{2}{p{0.3\linewidth}}{
        \cellcolor[rgb]{0.9,0.9,0.9}{
            \makecell[{{p{\linewidth}}}]{
                \texttt{\tiny{[GM$|$GM]}}
                \texttt{mücadele} \\
            }
        }
    }
    & & \\ \\

    \theutterance \stepcounter{utterance}  
    & & & \multicolumn{2}{p{0.3\linewidth}}{
        \cellcolor[rgb]{0.9,0.9,0.9}{
            \makecell[{{p{\linewidth}}}]{
                \texttt{\tiny{[GM$|$GM]}}
                ['devrim', 'savaş']
            }
        }
    }
    & & \\ \\

    \theutterance \stepcounter{utterance}  
    & & & \multicolumn{2}{p{0.3\linewidth}}{
        \cellcolor[rgb]{0.9,0.9,0.9}{
            \makecell[{{p{\linewidth}}}]{
                \texttt{\tiny{[GM$|$GM]}}
                {'word': 'devrim', 'assignment': 'team'}
            }
        }
    }
    & & \\ \\

    \theutterance \stepcounter{utterance}  
    & & & \multicolumn{2}{p{0.3\linewidth}}{
        \cellcolor[rgb]{0.9,0.9,0.9}{
            \makecell[{{p{\linewidth}}}]{
                \texttt{\tiny{[GM$|$GM]}}
                {'word': 'savaş', 'assignment': 'team'}
            }
        }
    }
    & & \\ \\

    \theutterance \stepcounter{utterance}  
    & & \multicolumn{4}{p{0.6\linewidth}}{
        \cellcolor[rgb]{0.9,0.9,0.9}{
            \makecell[{{p{\linewidth}}}]{
                \texttt{\tiny{[GM$\rangle$P2]}}
                \texttt{Aşağıdaki listeden, verilen sözcük "mücadele" ile en iyi ilgili veya en yakından bağlantılı, virgülle ayrılmış en fazla 2 sözcükten oluşan bir liste verin. Tahmininizi/tahmini(n)i her zaman TAHMİN: başlayın ve yanıtınıza başka hiçbir metin eklemeyin.} \\
\\ 
\texttt{stadyum, ayı, bıçak, kulüp, çapraz, öğrenci, bal, bilim adamı, ardıç kuşu, tokyo, kaçakçı, olympos, sıra, tahvil, antarktika, tarih, berlin, nal, savaş, süper kahraman, reçel, posta, taç, devrim, mikroskop} \\
            }
        }
    }
    & \\ \\

    \theutterance \stepcounter{utterance}  
    & & & \multicolumn{4}{p{0.6\linewidth}}{
        \cellcolor[rgb]{0.9,0.9,0.9}{
            \makecell[{{p{\linewidth}}}]{
                \texttt{\tiny{[GM$\langle$P2]}}
                \texttt{TAHMİN: savaş, devrim} \\
            }
        }
    }
     \\ \\

    \theutterance \stepcounter{utterance}  
    & & & \multicolumn{2}{p{0.3\linewidth}}{
        \cellcolor[rgb]{0.9,0.9,0.9}{
            \makecell[{{p{\linewidth}}}]{
                \texttt{\tiny{[GM$|$GM]}}
                ['savaş', 'devrim']
            }
        }
    }
    & & \\ \\

    \theutterance \stepcounter{utterance}  
    & & & \multicolumn{2}{p{0.3\linewidth}}{
        \cellcolor[rgb]{0.9,0.9,0.9}{
            \makecell[{{p{\linewidth}}}]{
                \texttt{\tiny{[GM$|$GM]}}
                {'word': 'savaş', 'assignment': 'team'}
            }
        }
    }
    & & \\ \\

    \theutterance \stepcounter{utterance}  
    & & & \multicolumn{2}{p{0.3\linewidth}}{
        \cellcolor[rgb]{0.9,0.9,0.9}{
            \makecell[{{p{\linewidth}}}]{
                \texttt{\tiny{[GM$|$GM]}}
                {'word': 'savaş', 'assignment': 'team'}
            }
        }
    }
    & & \\ \\

    \theutterance \stepcounter{utterance}  
    & & & \multicolumn{2}{p{0.3\linewidth}}{
        \cellcolor[rgb]{0.9,0.9,0.9}{
            \makecell[{{p{\linewidth}}}]{
                \texttt{\tiny{[GM$|$GM]}}
                {'word': 'devrim', 'assignment': 'team'}
            }
        }
    }
    & & \\ \\

    \theutterance \stepcounter{utterance}  
    & & & \multicolumn{2}{p{0.3\linewidth}}{
        \cellcolor[rgb]{0.9,0.9,0.9}{
            \makecell[{{p{\linewidth}}}]{
                \texttt{\tiny{[GM$|$GM]}}
                {'word': 'devrim', 'assignment': 'team'}
            }
        }
    }
    & & \\ \\

    \theutterance \stepcounter{utterance}  
    & & & \multicolumn{2}{p{0.3\linewidth}}{
        \cellcolor[rgb]{0.9,0.9,0.9}{
            \makecell[{{p{\linewidth}}}]{
                \texttt{\tiny{[GM$|$GM]}}
                {'word': 'bilim adamı', 'assignment': 'opponent'}
            }
        }
    }
    & & \\ \\

    \theutterance \stepcounter{utterance}  
    & & & \multicolumn{2}{p{0.3\linewidth}}{
        \cellcolor[rgb]{0.9,0.9,0.9}{
            \makecell[{{p{\linewidth}}}]{
                \texttt{\tiny{[GM$|$GM]}}
                {'hidden': {'team': ['kulüp', 'berlin', 'antarktika', 'stadyum', 'öğrenci', 'mikroskop', 'bal'], 'innocent': ['tokyo', 'posta', 'taç', 'süper kahraman', 'ardıç kuşu', 'nal', 'çapraz'], 'opponent': ['sıra', 'ayı', 'tahvil', 'bıçak', 'tarih', 'kaçakçı', 'olympos'], 'assassin': ['reçel']}, 'revealed': {'team': {'team': ['savaş', 'devrim'], 'innocent': [], 'opponent': [], 'assassin': []}, 'opponent': {'team': [], 'innocent': [], 'opponent': ['bilim adamı'], 'assassin': []}}}
            }
        }
    }
    & & \\ \\

    \theutterance \stepcounter{utterance}  
    & & \multicolumn{4}{p{0.6\linewidth}}{
        \cellcolor[rgb]{0.9,0.9,0.9}{
            \makecell[{{p{\linewidth}}}]{
                \texttt{\tiny{[P1$\langle$GM]}}
                \texttt{savaş, devrim kelimeleri doğru tahmin edildi. Takım arkadaşının sırası orada bitti.} \\
\texttt{Şimdi, kalan takım kelimelerinden bazılarıyla ilgili başka bir ipucu ve ilgili hedef kelimelerin bir listesini ver. İpucuna 'İPUCU: ' ile başlamayı, yeni bir satıra geçmeyi ve virgülle ayrılmış hedef kelime listene 'HEDEFLER: ' ile başlamayı unutma. Dikkat: önceki isteklere kıyasla listelerden bazı kelimeler çıkarılmış olabilir.} \\
\\ 
\texttt{Kalan takım kelimelerin: kulüp, berlin, antarktika, stadyum, öğrenci, mikroskop, bal.} \\
\texttt{Rakibin için kalan kelimeler: sıra, ayı, tahvil, bilim adamı, bıçak, tarih, kaçakçı, olympos.} \\
\texttt{Kalan dikkat dağıtıcı kelimeler: tokyo, posta, taç, süper kahraman, ardıç kuşu, nal, çapraz.} \\
\texttt{Kalan öldüren kelimeler: reçel.} \\
            }
        }
    }
    & \\ \\

    \theutterance \stepcounter{utterance}  
    & \multicolumn{4}{p{0.6\linewidth}}{
        \cellcolor[rgb]{0.9,0.9,0.9}{
            \makecell[{{p{\linewidth}}}]{
                \texttt{\tiny{[P1$\rangle$GM]}}
                \texttt{İPUCU: spor} \\
\texttt{HEDEFLER: kulüp, stadyum} \\
            }
        }
    }
    & & \\ \\

    \theutterance \stepcounter{utterance}  
    & & & \multicolumn{2}{p{0.3\linewidth}}{
        \cellcolor[rgb]{0.9,0.9,0.9}{
            \makecell[{{p{\linewidth}}}]{
                \texttt{\tiny{[GM$|$GM]}}
                \texttt{spor} \\
            }
        }
    }
    & & \\ \\

    \theutterance \stepcounter{utterance}  
    & & & \multicolumn{2}{p{0.3\linewidth}}{
        \cellcolor[rgb]{0.9,0.9,0.9}{
            \makecell[{{p{\linewidth}}}]{
                \texttt{\tiny{[GM$|$GM]}}
                ['kulüp', 'stadyum']
            }
        }
    }
    & & \\ \\

    \theutterance \stepcounter{utterance}  
    & & & \multicolumn{2}{p{0.3\linewidth}}{
        \cellcolor[rgb]{0.9,0.9,0.9}{
            \makecell[{{p{\linewidth}}}]{
                \texttt{\tiny{[GM$|$GM]}}
                {'word': 'kulüp', 'assignment': 'team'}
            }
        }
    }
    & & \\ \\

    \theutterance \stepcounter{utterance}  
    & & & \multicolumn{2}{p{0.3\linewidth}}{
        \cellcolor[rgb]{0.9,0.9,0.9}{
            \makecell[{{p{\linewidth}}}]{
                \texttt{\tiny{[GM$|$GM]}}
                {'word': 'stadyum', 'assignment': 'team'}
            }
        }
    }
    & & \\ \\

    \theutterance \stepcounter{utterance}  
    & & \multicolumn{4}{p{0.6\linewidth}}{
        \cellcolor[rgb]{0.9,0.9,0.9}{
            \makecell[{{p{\linewidth}}}]{
                \texttt{\tiny{[GM$\rangle$P2]}}
                \texttt{savaş, devrim kelimeleri doğru tahmin edildi. Senin sıran orada bitti.} \\
\texttt{Şimdi, aşağıdaki kelime listesinden, 'spor' sözcüğüyle en iyi ilişkili veya en yakından bağlantılı, en az 1 ve en fazla 2 kelimeden oluşan başka bir virgülle ayrılmış liste sağlayın. Cevabınızı TAHMİN: başlamayı unutmayın. Not: bazı kelimeler, önceki isteklerle karşılaştırıldığında listeden kaldırılmış olabilir.} \\
\\ 
\texttt{stadyum, ayı, bıçak, kulüp, çapraz, öğrenci, bal, ardıç kuşu, tokyo, kaçakçı, olympos, sıra, tahvil, antarktika, tarih, berlin, nal, süper kahraman, reçel, posta, taç, mikroskop} \\
            }
        }
    }
    & \\ \\

    \theutterance \stepcounter{utterance}  
    & & & \multicolumn{4}{p{0.6\linewidth}}{
        \cellcolor[rgb]{0.9,0.9,0.9}{
            \makecell[{{p{\linewidth}}}]{
                \texttt{\tiny{[GM$\langle$P2]}}
                \texttt{TAHMİN: stadyum, kulüp} \\
            }
        }
    }
     \\ \\

    \theutterance \stepcounter{utterance}  
    & & & \multicolumn{2}{p{0.3\linewidth}}{
        \cellcolor[rgb]{0.9,0.9,0.9}{
            \makecell[{{p{\linewidth}}}]{
                \texttt{\tiny{[GM$|$GM]}}
                ['stadyum', 'kulüp']
            }
        }
    }
    & & \\ \\

    \theutterance \stepcounter{utterance}  
    & & & \multicolumn{2}{p{0.3\linewidth}}{
        \cellcolor[rgb]{0.9,0.9,0.9}{
            \makecell[{{p{\linewidth}}}]{
                \texttt{\tiny{[GM$|$GM]}}
                {'word': 'stadyum', 'assignment': 'team'}
            }
        }
    }
    & & \\ \\

    \theutterance \stepcounter{utterance}  
    & & & \multicolumn{2}{p{0.3\linewidth}}{
        \cellcolor[rgb]{0.9,0.9,0.9}{
            \makecell[{{p{\linewidth}}}]{
                \texttt{\tiny{[GM$|$GM]}}
                {'word': 'stadyum', 'assignment': 'team'}
            }
        }
    }
    & & \\ \\

    \theutterance \stepcounter{utterance}  
    & & & \multicolumn{2}{p{0.3\linewidth}}{
        \cellcolor[rgb]{0.9,0.9,0.9}{
            \makecell[{{p{\linewidth}}}]{
                \texttt{\tiny{[GM$|$GM]}}
                {'word': 'kulüp', 'assignment': 'team'}
            }
        }
    }
    & & \\ \\

    \theutterance \stepcounter{utterance}  
    & & & \multicolumn{2}{p{0.3\linewidth}}{
        \cellcolor[rgb]{0.9,0.9,0.9}{
            \makecell[{{p{\linewidth}}}]{
                \texttt{\tiny{[GM$|$GM]}}
                {'word': 'kulüp', 'assignment': 'team'}
            }
        }
    }
    & & \\ \\

    \theutterance \stepcounter{utterance}  
    & & & \multicolumn{2}{p{0.3\linewidth}}{
        \cellcolor[rgb]{0.9,0.9,0.9}{
            \makecell[{{p{\linewidth}}}]{
                \texttt{\tiny{[GM$|$GM]}}
                {'word': 'kaçakçı', 'assignment': 'opponent'}
            }
        }
    }
    & & \\ \\

    \theutterance \stepcounter{utterance}  
    & & & \multicolumn{2}{p{0.3\linewidth}}{
        \cellcolor[rgb]{0.9,0.9,0.9}{
            \makecell[{{p{\linewidth}}}]{
                \texttt{\tiny{[GM$|$GM]}}
                {'hidden': {'team': ['berlin', 'antarktika', 'öğrenci', 'mikroskop', 'bal'], 'innocent': ['tokyo', 'posta', 'taç', 'süper kahraman', 'ardıç kuşu', 'nal', 'çapraz'], 'opponent': ['sıra', 'ayı', 'tahvil', 'bıçak', 'tarih', 'olympos'], 'assassin': ['reçel']}, 'revealed': {'team': {'team': ['savaş', 'devrim', 'stadyum', 'kulüp'], 'innocent': [], 'opponent': [], 'assassin': []}, 'opponent': {'team': [], 'innocent': [], 'opponent': ['bilim adamı', 'kaçakçı'], 'assassin': []}}}
            }
        }
    }
    & & \\ \\

    \theutterance \stepcounter{utterance}  
    & & \multicolumn{4}{p{0.6\linewidth}}{
        \cellcolor[rgb]{0.9,0.9,0.9}{
            \makecell[{{p{\linewidth}}}]{
                \texttt{\tiny{[P1$\langle$GM]}}
                \texttt{stadyum, kulüp kelimeleri doğru tahmin edildi. Takım arkadaşının sırası orada bitti.} \\
\texttt{Şimdi, kalan takım kelimelerinden bazılarıyla ilgili başka bir ipucu ve ilgili hedef kelimelerin bir listesini ver. İpucuna 'İPUCU: ' ile başlamayı, yeni bir satıra geçmeyi ve virgülle ayrılmış hedef kelime listene 'HEDEFLER: ' ile başlamayı unutma. Dikkat: önceki isteklere kıyasla listelerden bazı kelimeler çıkarılmış olabilir.} \\
\\ 
\texttt{Kalan takım kelimelerin: berlin, antarktika, öğrenci, mikroskop, bal.} \\
\texttt{Rakibin için kalan kelimeler: sıra, ayı, tahvil, bıçak, tarih, kaçakçı, olympos.} \\
\texttt{Kalan dikkat dağıtıcı kelimeler: tokyo, posta, taç, süper kahraman, ardıç kuşu, nal, çapraz.} \\
\texttt{Kalan öldüren kelimeler: reçel.} \\
            }
        }
    }
    & \\ \\

    \theutterance \stepcounter{utterance}  
    & \multicolumn{4}{p{0.6\linewidth}}{
        \cellcolor[rgb]{0.9,0.9,0.9}{
            \makecell[{{p{\linewidth}}}]{
                \texttt{\tiny{[P1$\rangle$GM]}}
                \texttt{İPUCU: laboratuvar} \\
\texttt{HEDEFLER: mikroskop, öğrenci} \\
            }
        }
    }
    & & \\ \\

    \theutterance \stepcounter{utterance}  
    & & & \multicolumn{2}{p{0.3\linewidth}}{
        \cellcolor[rgb]{0.9,0.9,0.9}{
            \makecell[{{p{\linewidth}}}]{
                \texttt{\tiny{[GM$|$GM]}}
                \texttt{laboratuvar} \\
            }
        }
    }
    & & \\ \\

    \theutterance \stepcounter{utterance}  
    & & & \multicolumn{2}{p{0.3\linewidth}}{
        \cellcolor[rgb]{0.9,0.9,0.9}{
            \makecell[{{p{\linewidth}}}]{
                \texttt{\tiny{[GM$|$GM]}}
                ['mikroskop', 'öğrenci']
            }
        }
    }
    & & \\ \\

    \theutterance \stepcounter{utterance}  
    & & & \multicolumn{2}{p{0.3\linewidth}}{
        \cellcolor[rgb]{0.9,0.9,0.9}{
            \makecell[{{p{\linewidth}}}]{
                \texttt{\tiny{[GM$|$GM]}}
                {'word': 'mikroskop', 'assignment': 'team'}
            }
        }
    }
    & & \\ \\

    \theutterance \stepcounter{utterance}  
    & & & \multicolumn{2}{p{0.3\linewidth}}{
        \cellcolor[rgb]{0.9,0.9,0.9}{
            \makecell[{{p{\linewidth}}}]{
                \texttt{\tiny{[GM$|$GM]}}
                {'word': 'öğrenci', 'assignment': 'team'}
            }
        }
    }
    & & \\ \\

    \theutterance \stepcounter{utterance}  
    & & \multicolumn{4}{p{0.6\linewidth}}{
        \cellcolor[rgb]{0.9,0.9,0.9}{
            \makecell[{{p{\linewidth}}}]{
                \texttt{\tiny{[GM$\rangle$P2]}}
                \texttt{stadyum, kulüp kelimeleri doğru tahmin edildi. Senin sıran orada bitti.} \\
\texttt{Şimdi, aşağıdaki kelime listesinden, 'laboratuvar' sözcüğüyle en iyi ilişkili veya en yakından bağlantılı, en az 1 ve en fazla 2 kelimeden oluşan başka bir virgülle ayrılmış liste sağlayın. Cevabınızı TAHMİN: başlamayı unutmayın. Not: bazı kelimeler, önceki isteklerle karşılaştırıldığında listeden kaldırılmış olabilir.} \\
\\ 
\texttt{ayı, bıçak, çapraz, öğrenci, bal, ardıç kuşu, tokyo, olympos, sıra, tahvil, antarktika, tarih, berlin, nal, süper kahraman, reçel, posta, taç, mikroskop} \\
            }
        }
    }
    & \\ \\

    \theutterance \stepcounter{utterance}  
    & & & \multicolumn{4}{p{0.6\linewidth}}{
        \cellcolor[rgb]{0.9,0.9,0.9}{
            \makecell[{{p{\linewidth}}}]{
                \texttt{\tiny{[GM$\langle$P2]}}
                \texttt{TAHMİN: mikroskop, öğrenci} \\
            }
        }
    }
     \\ \\

    \theutterance \stepcounter{utterance}  
    & & & \multicolumn{2}{p{0.3\linewidth}}{
        \cellcolor[rgb]{0.9,0.9,0.9}{
            \makecell[{{p{\linewidth}}}]{
                \texttt{\tiny{[GM$|$GM]}}
                ['mikroskop', 'öğrenci']
            }
        }
    }
    & & \\ \\

    \theutterance \stepcounter{utterance}  
    & & & \multicolumn{2}{p{0.3\linewidth}}{
        \cellcolor[rgb]{0.9,0.9,0.9}{
            \makecell[{{p{\linewidth}}}]{
                \texttt{\tiny{[GM$|$GM]}}
                {'word': 'mikroskop', 'assignment': 'team'}
            }
        }
    }
    & & \\ \\

    \theutterance \stepcounter{utterance}  
    & & & \multicolumn{2}{p{0.3\linewidth}}{
        \cellcolor[rgb]{0.9,0.9,0.9}{
            \makecell[{{p{\linewidth}}}]{
                \texttt{\tiny{[GM$|$GM]}}
                {'word': 'mikroskop', 'assignment': 'team'}
            }
        }
    }
    & & \\ \\

    \theutterance \stepcounter{utterance}  
    & & & \multicolumn{2}{p{0.3\linewidth}}{
        \cellcolor[rgb]{0.9,0.9,0.9}{
            \makecell[{{p{\linewidth}}}]{
                \texttt{\tiny{[GM$|$GM]}}
                {'word': 'öğrenci', 'assignment': 'team'}
            }
        }
    }
    & & \\ \\

    \theutterance \stepcounter{utterance}  
    & & & \multicolumn{2}{p{0.3\linewidth}}{
        \cellcolor[rgb]{0.9,0.9,0.9}{
            \makecell[{{p{\linewidth}}}]{
                \texttt{\tiny{[GM$|$GM]}}
                {'word': 'öğrenci', 'assignment': 'team'}
            }
        }
    }
    & & \\ \\

    \theutterance \stepcounter{utterance}  
    & & & \multicolumn{2}{p{0.3\linewidth}}{
        \cellcolor[rgb]{0.9,0.9,0.9}{
            \makecell[{{p{\linewidth}}}]{
                \texttt{\tiny{[GM$|$GM]}}
                {'word': 'tahvil', 'assignment': 'opponent'}
            }
        }
    }
    & & \\ \\

    \theutterance \stepcounter{utterance}  
    & & & \multicolumn{2}{p{0.3\linewidth}}{
        \cellcolor[rgb]{0.9,0.9,0.9}{
            \makecell[{{p{\linewidth}}}]{
                \texttt{\tiny{[GM$|$GM]}}
                {'hidden': {'team': ['berlin', 'antarktika', 'bal'], 'innocent': ['tokyo', 'posta', 'taç', 'süper kahraman', 'ardıç kuşu', 'nal', 'çapraz'], 'opponent': ['sıra', 'ayı', 'bıçak', 'tarih', 'olympos'], 'assassin': ['reçel']}, 'revealed': {'team': {'team': ['savaş', 'devrim', 'stadyum', 'kulüp', 'mikroskop', 'öğrenci'], 'innocent': [], 'opponent': [], 'assassin': []}, 'opponent': {'team': [], 'innocent': [], 'opponent': ['bilim adamı', 'kaçakçı', 'tahvil'], 'assassin': []}}}
            }
        }
    }
    & & \\ \\

    \theutterance \stepcounter{utterance}  
    & & \multicolumn{4}{p{0.6\linewidth}}{
        \cellcolor[rgb]{0.9,0.9,0.9}{
            \makecell[{{p{\linewidth}}}]{
                \texttt{\tiny{[P1$\langle$GM]}}
                \texttt{mikroskop, öğrenci kelimeleri doğru tahmin edildi. Takım arkadaşının sırası orada bitti.} \\
\texttt{Şimdi, kalan takım kelimelerinden bazılarıyla ilgili başka bir ipucu ve ilgili hedef kelimelerin bir listesini ver. İpucuna 'İPUCU: ' ile başlamayı, yeni bir satıra geçmeyi ve virgülle ayrılmış hedef kelime listene 'HEDEFLER: ' ile başlamayı unutma. Dikkat: önceki isteklere kıyasla listelerden bazı kelimeler çıkarılmış olabilir.} \\
\\ 
\texttt{Kalan takım kelimelerin: berlin, antarktika, bal.} \\
\texttt{Rakibin için kalan kelimeler: sıra, ayı, tahvil, bıçak, tarih, olympos.} \\
\texttt{Kalan dikkat dağıtıcı kelimeler: tokyo, posta, taç, süper kahraman, ardıç kuşu, nal, çapraz.} \\
\texttt{Kalan öldüren kelimeler: reçel.} \\
            }
        }
    }
    & \\ \\

    \theutterance \stepcounter{utterance}  
    & \multicolumn{4}{p{0.6\linewidth}}{
        \cellcolor[rgb]{0.9,0.9,0.9}{
            \makecell[{{p{\linewidth}}}]{
                \texttt{\tiny{[P1$\rangle$GM]}}
                \texttt{İPUCU: soğuk} \\
\texttt{HEDEFLER: berlin, antarktika} \\
            }
        }
    }
    & & \\ \\

    \theutterance \stepcounter{utterance}  
    & & & \multicolumn{2}{p{0.3\linewidth}}{
        \cellcolor[rgb]{0.9,0.9,0.9}{
            \makecell[{{p{\linewidth}}}]{
                \texttt{\tiny{[GM$|$GM]}}
                \texttt{soğuk} \\
            }
        }
    }
    & & \\ \\

    \theutterance \stepcounter{utterance}  
    & & & \multicolumn{2}{p{0.3\linewidth}}{
        \cellcolor[rgb]{0.9,0.9,0.9}{
            \makecell[{{p{\linewidth}}}]{
                \texttt{\tiny{[GM$|$GM]}}
                ['berlin', 'antarktika']
            }
        }
    }
    & & \\ \\

    \theutterance \stepcounter{utterance}  
    & & & \multicolumn{2}{p{0.3\linewidth}}{
        \cellcolor[rgb]{0.9,0.9,0.9}{
            \makecell[{{p{\linewidth}}}]{
                \texttt{\tiny{[GM$|$GM]}}
                {'word': 'berlin', 'assignment': 'team'}
            }
        }
    }
    & & \\ \\

    \theutterance \stepcounter{utterance}  
    & & & \multicolumn{2}{p{0.3\linewidth}}{
        \cellcolor[rgb]{0.9,0.9,0.9}{
            \makecell[{{p{\linewidth}}}]{
                \texttt{\tiny{[GM$|$GM]}}
                {'word': 'antarktika', 'assignment': 'team'}
            }
        }
    }
    & & \\ \\

    \theutterance \stepcounter{utterance}  
    & & \multicolumn{4}{p{0.6\linewidth}}{
        \cellcolor[rgb]{0.9,0.9,0.9}{
            \makecell[{{p{\linewidth}}}]{
                \texttt{\tiny{[GM$\rangle$P2]}}
                \texttt{mikroskop, öğrenci kelimeleri doğru tahmin edildi. Senin sıran orada bitti.} \\
\texttt{Şimdi, aşağıdaki kelime listesinden, 'soğuk' sözcüğüyle en iyi ilişkili veya en yakından bağlantılı, en az 1 ve en fazla 2 kelimeden oluşan başka bir virgülle ayrılmış liste sağlayın. Cevabınızı TAHMİN: başlamayı unutmayın. Not: bazı kelimeler, önceki isteklerle karşılaştırıldığında listeden kaldırılmış olabilir.} \\
\\ 
\texttt{ayı, bıçak, çapraz, bal, ardıç kuşu, tokyo, olympos, sıra, antarktika, tarih, berlin, nal, süper kahraman, reçel, posta, taç} \\
            }
        }
    }
    & \\ \\

    \theutterance \stepcounter{utterance}  
    & & & \multicolumn{4}{p{0.6\linewidth}}{
        \cellcolor[rgb]{0.9,0.9,0.9}{
            \makecell[{{p{\linewidth}}}]{
                \texttt{\tiny{[GM$\langle$P2]}}
                \texttt{TAHMİN: antarktika, ayı} \\
            }
        }
    }
     \\ \\

    \theutterance \stepcounter{utterance}  
    & & & \multicolumn{2}{p{0.3\linewidth}}{
        \cellcolor[rgb]{0.9,0.9,0.9}{
            \makecell[{{p{\linewidth}}}]{
                \texttt{\tiny{[GM$|$GM]}}
                ['antarktika', 'ayı']
            }
        }
    }
    & & \\ \\

    \theutterance \stepcounter{utterance}  
    & & & \multicolumn{2}{p{0.3\linewidth}}{
        \cellcolor[rgb]{0.9,0.9,0.9}{
            \makecell[{{p{\linewidth}}}]{
                \texttt{\tiny{[GM$|$GM]}}
                {'word': 'antarktika', 'assignment': 'team'}
            }
        }
    }
    & & \\ \\

    \theutterance \stepcounter{utterance}  
    & & & \multicolumn{2}{p{0.3\linewidth}}{
        \cellcolor[rgb]{0.9,0.9,0.9}{
            \makecell[{{p{\linewidth}}}]{
                \texttt{\tiny{[GM$|$GM]}}
                {'word': 'antarktika', 'assignment': 'team'}
            }
        }
    }
    & & \\ \\

    \theutterance \stepcounter{utterance}  
    & & & \multicolumn{2}{p{0.3\linewidth}}{
        \cellcolor[rgb]{0.9,0.9,0.9}{
            \makecell[{{p{\linewidth}}}]{
                \texttt{\tiny{[GM$|$GM]}}
                {'word': 'ayı', 'assignment': 'opponent'}
            }
        }
    }
    & & \\ \\

    \theutterance \stepcounter{utterance}  
    & & & \multicolumn{2}{p{0.3\linewidth}}{
        \cellcolor[rgb]{0.9,0.9,0.9}{
            \makecell[{{p{\linewidth}}}]{
                \texttt{\tiny{[GM$|$GM]}}
                \texttt{ayı} \\
            }
        }
    }
    & & \\ \\

    \theutterance \stepcounter{utterance}  
    & & & \multicolumn{2}{p{0.3\linewidth}}{
        \cellcolor[rgb]{0.9,0.9,0.9}{
            \makecell[{{p{\linewidth}}}]{
                \texttt{\tiny{[GM$|$GM]}}
                {'word': 'sıra', 'assignment': 'opponent'}
            }
        }
    }
    & & \\ \\

    \theutterance \stepcounter{utterance}  
    & & & \multicolumn{2}{p{0.3\linewidth}}{
        \cellcolor[rgb]{0.9,0.9,0.9}{
            \makecell[{{p{\linewidth}}}]{
                \texttt{\tiny{[GM$|$GM]}}
                {'hidden': {'team': ['berlin', 'bal'], 'innocent': ['tokyo', 'posta', 'taç', 'süper kahraman', 'ardıç kuşu', 'nal', 'çapraz'], 'opponent': ['bıçak', 'tarih', 'olympos'], 'assassin': ['reçel']}, 'revealed': {'team': {'team': ['savaş', 'devrim', 'stadyum', 'kulüp', 'mikroskop', 'öğrenci', 'antarktika'], 'innocent': [], 'opponent': ['ayı'], 'assassin': []}, 'opponent': {'team': [], 'innocent': [], 'opponent': ['bilim adamı', 'kaçakçı', 'tahvil', 'sıra'], 'assassin': []}}}
            }
        }
    }
    & & \\ \\

    \theutterance \stepcounter{utterance}  
    & & \multicolumn{4}{p{0.6\linewidth}}{
        \cellcolor[rgb]{0.9,0.9,0.9}{
            \makecell[{{p{\linewidth}}}]{
                \texttt{\tiny{[P1$\langle$GM]}}
                \texttt{antarktika kelimesi doğru tahmin edildi. ayı kelimesi tahmin edildi, ancak bu bir rakip kelimedir. Takım arkadaşının sırası orada bitti.} \\
\texttt{Şimdi, kalan takım kelimelerinden bazılarıyla ilgili başka bir ipucu ve ilgili hedef kelimelerin bir listesini ver. İpucuna 'İPUCU: ' ile başlamayı, yeni bir satıra geçmeyi ve virgülle ayrılmış hedef kelime listene 'HEDEFLER: ' ile başlamayı unutma. Dikkat: önceki isteklere kıyasla listelerden bazı kelimeler çıkarılmış olabilir.} \\
\\ 
\texttt{Kalan takım kelimelerin: berlin, bal.} \\
\texttt{Rakibin için kalan kelimeler: sıra, bıçak, tarih, olympos.} \\
\texttt{Kalan dikkat dağıtıcı kelimeler: tokyo, posta, taç, süper kahraman, ardıç kuşu, nal, çapraz.} \\
\texttt{Kalan öldüren kelimeler: reçel.} \\
            }
        }
    }
    & \\ \\

    \theutterance \stepcounter{utterance}  
    & \multicolumn{4}{p{0.6\linewidth}}{
        \cellcolor[rgb]{0.9,0.9,0.9}{
            \makecell[{{p{\linewidth}}}]{
                \texttt{\tiny{[P1$\rangle$GM]}}
                \texttt{İPUCU: arı} \\
\texttt{HEDEFLER: bal} \\
            }
        }
    }
    & & \\ \\

    \theutterance \stepcounter{utterance}  
    & & & \multicolumn{2}{p{0.3\linewidth}}{
        \cellcolor[rgb]{0.9,0.9,0.9}{
            \makecell[{{p{\linewidth}}}]{
                \texttt{\tiny{[GM$|$GM]}}
                \texttt{arı} \\
            }
        }
    }
    & & \\ \\

    \theutterance \stepcounter{utterance}  
    & & & \multicolumn{2}{p{0.3\linewidth}}{
        \cellcolor[rgb]{0.9,0.9,0.9}{
            \makecell[{{p{\linewidth}}}]{
                \texttt{\tiny{[GM$|$GM]}}
                ['bal']
            }
        }
    }
    & & \\ \\

    \theutterance \stepcounter{utterance}  
    & & & \multicolumn{2}{p{0.3\linewidth}}{
        \cellcolor[rgb]{0.9,0.9,0.9}{
            \makecell[{{p{\linewidth}}}]{
                \texttt{\tiny{[GM$|$GM]}}
                {'word': 'bal', 'assignment': 'team'}
            }
        }
    }
    & & \\ \\

    \theutterance \stepcounter{utterance}  
    & & \multicolumn{4}{p{0.6\linewidth}}{
        \cellcolor[rgb]{0.9,0.9,0.9}{
            \makecell[{{p{\linewidth}}}]{
                \texttt{\tiny{[GM$\rangle$P2]}}
                \texttt{antarktika kelimesi doğru tahmin edildi. ayı kelimesi tahmin edildi, ancak bu bir rakip kelimedir. Senin sıran orada bitti.} \\
\texttt{Şimdi, aşağıdaki kelime listesinden, 'arı' sözcüğüyle en iyi ilişkili veya en yakından bağlantılı, en az 1 ve en fazla 1 kelimeden oluşan başka bir virgülle ayrılmış liste sağlayın. Cevabınızı TAHMİN: başlamayı unutmayın. Not: bazı kelimeler, önceki isteklerle karşılaştırıldığında listeden kaldırılmış olabilir.} \\
\\ 
\texttt{bıçak, çapraz, bal, ardıç kuşu, tokyo, olympos, tarih, berlin, nal, süper kahraman, reçel, posta, taç} \\
            }
        }
    }
    & \\ \\

    \theutterance \stepcounter{utterance}  
    & & & \multicolumn{4}{p{0.6\linewidth}}{
        \cellcolor[rgb]{0.9,0.9,0.9}{
            \makecell[{{p{\linewidth}}}]{
                \texttt{\tiny{[GM$\langle$P2]}}
                \texttt{TAHMİN: bal} \\
            }
        }
    }
     \\ \\

    \theutterance \stepcounter{utterance}  
    & & & \multicolumn{2}{p{0.3\linewidth}}{
        \cellcolor[rgb]{0.9,0.9,0.9}{
            \makecell[{{p{\linewidth}}}]{
                \texttt{\tiny{[GM$|$GM]}}
                ['bal']
            }
        }
    }
    & & \\ \\

    \theutterance \stepcounter{utterance}  
    & & & \multicolumn{2}{p{0.3\linewidth}}{
        \cellcolor[rgb]{0.9,0.9,0.9}{
            \makecell[{{p{\linewidth}}}]{
                \texttt{\tiny{[GM$|$GM]}}
                {'word': 'bal', 'assignment': 'team'}
            }
        }
    }
    & & \\ \\

    \theutterance \stepcounter{utterance}  
    & & & \multicolumn{2}{p{0.3\linewidth}}{
        \cellcolor[rgb]{0.9,0.9,0.9}{
            \makecell[{{p{\linewidth}}}]{
                \texttt{\tiny{[GM$|$GM]}}
                {'word': 'bal', 'assignment': 'team'}
            }
        }
    }
    & & \\ \\

    \theutterance \stepcounter{utterance}  
    & & & \multicolumn{2}{p{0.3\linewidth}}{
        \cellcolor[rgb]{0.9,0.9,0.9}{
            \makecell[{{p{\linewidth}}}]{
                \texttt{\tiny{[GM$|$GM]}}
                {'word': 'bıçak', 'assignment': 'opponent'}
            }
        }
    }
    & & \\ \\

    \theutterance \stepcounter{utterance}  
    & & & \multicolumn{2}{p{0.3\linewidth}}{
        \cellcolor[rgb]{0.9,0.9,0.9}{
            \makecell[{{p{\linewidth}}}]{
                \texttt{\tiny{[GM$|$GM]}}
                {'hidden': {'team': ['berlin'], 'innocent': ['tokyo', 'posta', 'taç', 'süper kahraman', 'ardıç kuşu', 'nal', 'çapraz'], 'opponent': ['tarih', 'olympos'], 'assassin': ['reçel']}, 'revealed': {'team': {'team': ['savaş', 'devrim', 'stadyum', 'kulüp', 'mikroskop', 'öğrenci', 'antarktika', 'bal'], 'innocent': [], 'opponent': ['ayı'], 'assassin': []}, 'opponent': {'team': [], 'innocent': [], 'opponent': ['bilim adamı', 'kaçakçı', 'tahvil', 'sıra', 'bıçak'], 'assassin': []}}}
            }
        }
    }
    & & \\ \\

    \theutterance \stepcounter{utterance}  
    & & \multicolumn{4}{p{0.6\linewidth}}{
        \cellcolor[rgb]{0.9,0.9,0.9}{
            \makecell[{{p{\linewidth}}}]{
                \texttt{\tiny{[P1$\langle$GM]}}
                \texttt{bal kelimesi doğru tahmin edildi. Takım arkadaşının sırası orada bitti.} \\
\texttt{Şimdi, kalan takım kelimelerinden bazılarıyla ilgili başka bir ipucu ve ilgili hedef kelimelerin bir listesini ver. İpucuna 'İPUCU: ' ile başlamayı, yeni bir satıra geçmeyi ve virgülle ayrılmış hedef kelime listene 'HEDEFLER: ' ile başlamayı unutma. Dikkat: önceki isteklere kıyasla listelerden bazı kelimeler çıkarılmış olabilir.} \\
\\ 
\texttt{Kalan takım kelimelerin: berlin.} \\
\texttt{Rakibin için kalan kelimeler: bıçak, tarih, olympos.} \\
\texttt{Kalan dikkat dağıtıcı kelimeler: tokyo, posta, taç, süper kahraman, ardıç kuşu, nal, çapraz.} \\
\texttt{Kalan öldüren kelimeler: reçel.} \\
            }
        }
    }
    & \\ \\

    \theutterance \stepcounter{utterance}  
    & \multicolumn{4}{p{0.6\linewidth}}{
        \cellcolor[rgb]{0.9,0.9,0.9}{
            \makecell[{{p{\linewidth}}}]{
                \texttt{\tiny{[P1$\rangle$GM]}}
                \texttt{İPUCU: almanya} \\
\texttt{HEDEFLER: berlin} \\
            }
        }
    }
    & & \\ \\

    \theutterance \stepcounter{utterance}  
    & & & \multicolumn{2}{p{0.3\linewidth}}{
        \cellcolor[rgb]{0.9,0.9,0.9}{
            \makecell[{{p{\linewidth}}}]{
                \texttt{\tiny{[GM$|$GM]}}
                \texttt{almanya} \\
            }
        }
    }
    & & \\ \\

    \theutterance \stepcounter{utterance}  
    & & & \multicolumn{2}{p{0.3\linewidth}}{
        \cellcolor[rgb]{0.9,0.9,0.9}{
            \makecell[{{p{\linewidth}}}]{
                \texttt{\tiny{[GM$|$GM]}}
                ['berlin']
            }
        }
    }
    & & \\ \\

    \theutterance \stepcounter{utterance}  
    & & & \multicolumn{2}{p{0.3\linewidth}}{
        \cellcolor[rgb]{0.9,0.9,0.9}{
            \makecell[{{p{\linewidth}}}]{
                \texttt{\tiny{[GM$|$GM]}}
                {'word': 'berlin', 'assignment': 'team'}
            }
        }
    }
    & & \\ \\

    \theutterance \stepcounter{utterance}  
    & & \multicolumn{4}{p{0.6\linewidth}}{
        \cellcolor[rgb]{0.9,0.9,0.9}{
            \makecell[{{p{\linewidth}}}]{
                \texttt{\tiny{[GM$\rangle$P2]}}
                \texttt{bal kelimesi doğru tahmin edildi. Senin sıran orada bitti.} \\
\texttt{Şimdi, aşağıdaki kelime listesinden, 'almanya' sözcüğüyle en iyi ilişkili veya en yakından bağlantılı, en az 1 ve en fazla 1 kelimeden oluşan başka bir virgülle ayrılmış liste sağlayın. Cevabınızı TAHMİN: başlamayı unutmayın. Not: bazı kelimeler, önceki isteklerle karşılaştırıldığında listeden kaldırılmış olabilir.} \\
\\ 
\texttt{çapraz, ardıç kuşu, tokyo, olympos, tarih, berlin, nal, süper kahraman, reçel, posta, taç} \\
            }
        }
    }
    & \\ \\

    \theutterance \stepcounter{utterance}  
    & & & \multicolumn{4}{p{0.6\linewidth}}{
        \cellcolor[rgb]{0.9,0.9,0.9}{
            \makecell[{{p{\linewidth}}}]{
                \texttt{\tiny{[GM$\langle$P2]}}
                \texttt{TAHMİN: berlin} \\
            }
        }
    }
     \\ \\

    \theutterance \stepcounter{utterance}  
    & & & \multicolumn{2}{p{0.3\linewidth}}{
        \cellcolor[rgb]{0.9,0.9,0.9}{
            \makecell[{{p{\linewidth}}}]{
                \texttt{\tiny{[GM$|$GM]}}
                ['berlin']
            }
        }
    }
    & & \\ \\

    \theutterance \stepcounter{utterance}  
    & & & \multicolumn{2}{p{0.3\linewidth}}{
        \cellcolor[rgb]{0.9,0.9,0.9}{
            \makecell[{{p{\linewidth}}}]{
                \texttt{\tiny{[GM$|$GM]}}
                {'word': 'berlin', 'assignment': 'team'}
            }
        }
    }
    & & \\ \\

    \theutterance \stepcounter{utterance}  
    & & & \multicolumn{2}{p{0.3\linewidth}}{
        \cellcolor[rgb]{0.9,0.9,0.9}{
            \makecell[{{p{\linewidth}}}]{
                \texttt{\tiny{[GM$|$GM]}}
                {'word': 'berlin', 'assignment': 'team'}
            }
        }
    }
    & & \\ \\

    \theutterance \stepcounter{utterance}  
    & & & \multicolumn{2}{p{0.3\linewidth}}{
        \cellcolor[rgb]{0.9,0.9,0.9}{
            \makecell[{{p{\linewidth}}}]{
                \texttt{\tiny{[GM$|$GM]}}
                \texttt{team has won} \\
            }
        }
    }
    & & \\ \\

\end{supertabular}
}

\end{document}
